%% start of file `template.tex'.
%% Copyright 2006-2013 Xavier Danaux (xdanaux@gmail.com).
%
% This work may be distributed and/or modified under the
% conditions of the LaTeX Project Public License version 1.3c,
% available at http://www.latex-project.org/lppl/.

\documentclass[11pt,a4paper,sans]{moderncv}
\moderncvstyle{banking}
\moderncvcolor{blue}
\usepackage[utf8]{inputenc}
\usepackage[scale=0.75, margin=0.8in]{geometry}
\usepackage{import}
\usepackage{fontawesome}
\usepackage{footmisc}

\name{A\textcolor{gray}{lexandre}}{G\textcolor{gray}{baguidi} A\textcolor{gray}{ïsse}}
\title{CV}
\address{02 Place Aurélie Nemours, 75013, Paris, France}{}{}
\phone[mobile]{+33 642426980}
\email{aga@lrde.epita.fr}
\homepage{shumpaga.com}
\social[linkedin]{alexandregbaguidiaisse}
\social[github]{shumpaga}

\usepackage{background}
\usepackage{color}
\usetikzlibrary{calc}
\newsavebox\mybox
\sbox\mybox{%
\begin{tikzpicture}[remember picture]
    \fill[cyan] (current page.north west) rectangle ($(current page.south west) +(10mm,0)$);
\end{tikzpicture}%
}
\backgroundsetup{
angle=0,
scale=1,
opacity=1,
color=black,
contents={%
\begin{tikzpicture}[remember picture, overlay]
    \node[anchor=west] at  (current page.west) {\usebox{\mybox}};
\end{tikzpicture}%
}
}
\usepackage{amsmath}
\newcommand{\CC }{\emph{ContextCapture$^{\text{\textbf{TM}}}$ }}
\usepackage{enumitem}
\definecolor{cusblue}{rgb}{0.0, 0.4, 0.6}

\begin{document}
\makecvtitle
\vspace{-1cm}
\begin{center}
  Ingénieur logiciel à \textbf{Acute3D | Bentley Systems}
\end{center}

\section{Experience}
\footnotetext[1]{Spot est une bibliothèque C++14 de vérification formelle par $\omega$-automates et formules LTL: \url{spot.lrde.epita.fr}.}
\begin{itemize}
\item{
    \cventry{Sep 2018 - Aujourd'hui}{\textcolor{cusblue}{Ingénieur logiciel}}{Bentley Systems -- Reality Modeling}{Paris, France}{}{\vspace{-2.8pt}
      Travaille actuellement sur la parallélisation des tâches dans \CC: scheduler, multithreading.
    }
    \vspace{-0.7cm}
    \cventry{Feb 2018 - Aug 2018}{\textcolor{cusblue}{Stage de fin d'étude}}{}{}{}{\vspace{-2.8pt}
      Amélioration du support de nuage de point dans \CC: un algorithme de \textbf{localisation de tous les scanners} d'un nuage de point statique et un algorithme de \textbf{visibilité} qui attribue à chaque point le scanner qui le voit le mieux. Ces deux algorithmes sont en cours de \textbf{brevetage}.
    }
  }

\item{\cventry{Sep 2016 -- Jan 2018}{\textcolor{cusblue}{Etudiant chercheur}}{LRDE (Laboratoire de Recherche et Développement de l'EPITA)}{Paris, France}{}{\vspace{-2.8pt}Plusieurs contributions sur Spot\footnotemark[1]: Optimisation d'un algorithme de minimisation d'$\omega$-automates utilisant un solveur de statisfaisabilité booléenne, conversion des classes communes d'$\omega$-automates en automates de co-Büchi et enfin, vérification de la propriété de ``persistence'' sur les Formules LTL.}}

\item{\cventry{Aug 2016}{\textcolor{cusblue}{Assistant de l'enseignant}}{Ecole polytechnique -- HEC}{Paris, France}{}{\vspace{-2.8pt}Une équipe de 4 assistants aidant les étudiants à rapidement prendre en main les notions de programmation (Python) enseignées dans un cours de Big Data.}}

\end{itemize}


\section{Education}
\begin{itemize}
\setlength\itemsep{-0.2cm}
\item{\cventry{Sep 2013 -- Sep 2018}{Ingénieurie et Science Informatique}{EPITA}{Paris, France} {\textit{Diplôme d'ingénieur}}{}}
\item{\cventry{Jan 2015 -- Mai 2015}{Temps réel et systèmes embarqués}{Staffordshire University}{Stafford, England} {\textit{Semestre à l'étranger}}{}}
\end{itemize}


\section{Quelques projets de programmation:}
\setlength{\tabcolsep}{3pt}
\renewcommand{\arraystretch}{1.2}
\begin{tabular}{ll}
  \textbf{Python:} & Détection de multiples code barre dans une image (using OpenCV).\\
& Blog personnel: \url{shumpaga.com} (module Django).\\
\textbf{Scala:} & Implémentation d'une IA 'Tron': 107ème place sur 4000 participants mondiaux.\\
  \textbf{C++:} & Allocateur de mémoire distribué avec ordonnancement causal des évènements (OpenMPI).\\
& Raytracer temps réel CUDA.\\
& Vérificateur d'orthographe basé sur un arbre PATRICIA, contrainte de 3000 qps/core et 512 Mo.\\
\textbf{C:} & Un shell POSIX-compliant basé sur la SCL.\\
& Implémentation de malloc, realloc, calloc et free avec la méthode de 'first fit'.\\
& Logiciel de reconnaissance optique de caractère utilisant un réseau neuronnal prédictif.\\
\end{tabular}

\section{Compétences}
\setlength{\tabcolsep}{3pt}
\renewcommand{\arraystretch}{1.2}
\begin{tabular}{ll}
\textbf{Languages:} & Français (Natif), Russe (Natif), Anglais (TOEIC 850)\\
\textbf{Programmation:} & C, C++, Python, MATLAB, C\#, Scala, PostgreSQL, Java.\\
\textbf{OS:} & Linux / UNIX, Mac Os X, Windows.\\
\textbf{Intérêts:} & Enregistrement audio, mix, master, musique assistée par ordinateur, voyages, évènementiel.
\end{tabular}

\end{document}
